\scanpage{12}\mspage{11}
% ink opening line
Pour ceci, on note \deleted{\unsure{qu'une str}}\inserted{d'abord}

% ink Lemme, first four lines
\uline{Lemme} Une structure de \deleted{biordre} tronçon biordonnien est connue quand on connaît la relation en \deleted{$\{a,b,c\}$} trois \unsure{élémt|éléments} $a,b,c$
% ink relation R, first display
\[
\begin{aligned}
\left.\vphantom{\substack{X\\X}}\right|\!\left.\vphantom{\substack{X\\X}}\right|\!\left.\vphantom{\substack{X\\X}}\right|\, R(a,b,c)
&\qquad b\ \text{\uline{est strictement entre}}\\
&\qquad\qquad\text{\uline{$a$ et $c$}}
\end{aligned}
\]
(au sens que $a<b<\deleted{\unsure{e}}c$ \uline{ou} $a>b>c$)
% ink left, short diagonal note beside the three strokes
% slanted
\margin{relation sym. en $a,c$}
% ink left, NB, diagonal block
% slanted
\margin{\uline{NB} \deleted{\unsure{On a de}}\inserted{La conn. de} Cette relation $R$ implique la relation $\bar R$ non stricte \deleted{\unsure{($a=b=c$} \unsure{ou}}\ \deleted{$\unsure{a=b<c}$}\ $a\leq b\leq c$ ou $a\geq b\geq c$ \deleted{$\unsure{a\leq b\leq c}$} $(\unsure{\rightleftarrows|\Leftrightarrow})$ \deleted{$\{a,b,c\}$} $b$ entre $a$ et $c$ \inserted{\uline{strict}} ou $\{a,b,c\}$ est un drapeau de cardinal $\leq 2$}

% ink middle, Dém and équivalence
\uline{Dém} \deleted{Choisissons} On en déduit la relation (i.e. \inserted{la conn. de} $S_0(\mathcal L)$\deleted{\unsure{)}} i.e. de $\mathrm{Drap}(\mathcal L)$)
\[
\begin{aligned}
\deleted{\unsure{a}}\,a\text{ et }b\ \text{``disjoints''}\quad&
\text{i.e. distincts et comparables}\\
&\text{(i.e. }\{a,b\}\in\mathrm{Drap}_2(\mathcal L)\\
&\text{pour le biordre}
\end{aligned}
\]
% ink lower middle, formule ssi
par 
\[
\begin{aligned}
a\text{ et }b\ \text{disjts}\quad \text{ssi}\quad
\exists\,c\ \text{tel que }c\text{ entre}\\
a\text{ et }b
\end{aligned}
\]
(utilise divisibilité ! d\deleted{\unsure{es}}\unsure{u} \deleted{\unsure{drap. le ré}} biordre)

% ink bottom, Choisissons
\deleted{\illegible{1}} Choisissons $\unsure{\varepsilon|\mathcal E|E}\in S_0$ (existe, car \deleted{si $x,y\in$} par hyp $\exists\,x,y\in\mathcal L,\quad x\ne y$, donc $a\leq x,y\leq b$, et on a $a\leq b$ et $a\ne b$ \deleted{\unsure{)}} \deleted{\unsure{c|p}}car sinon on aurait $x=y$)
% ink lower left, four short oblique lines
% slanted
\margin{Il existe $=$ des \unsure{drap.} de long. arbitraire}

% pale vertical watermark
\rotatebox{90}{\textsc{\unsure{ra Strong Extra}}}
